% Simplified galvanic series, with the three pairings that matter in airframes.
\begin{tikzpicture}[>=Latex]
\tikzset{
  mt/.style={font=\small,anchor=west},
  tick/.style={black!60,thin},
  good/.style={draw=green!45!black,line width=1.1pt},
  bad/.style={draw=red!70!black,line width=1.1pt},
  gl/.style={font=\scriptsize,align=left,text width=5.0cm,anchor=west,text=green!35!black},
  bl/.style={font=\scriptsize,align=left,text width=5.0cm,anchor=west,text=red!65!black}
}
\def\W{0.85}

% ---- the strip ----
\shade[top color=red!55,bottom color=blue!55] (0,0) rectangle (\W,10.4);
\draw[black!60] (0,0) rectangle (\W,10.4);
\node[font=\scriptsize\bfseries,text=red!65!black,anchor=south,align=center] at (\W/2,10.55) {ANODIC\\[-1pt] \mdseries\itshape this one corrodes};
\node[font=\scriptsize\bfseries,text=blue!60!black,anchor=north,align=center] at (\W/2,-0.15) {CATHODIC\\[-1pt] \mdseries\itshape this one is protected};

% ---- members ----
\foreach \y/\n in {
  10.0/{Magnesium},
   9.2/{Zinc},
   8.5/{Aluminium 1100},
   8.0/{\textbf{Cadmium plating}},
   7.4/{Aluminium 2024, 7075},
   6.3/{Carbon and low-alloy steel},
   5.0/{300-series CRES \emph{(passive)}},
   4.2/{A286},
   3.3/{\textbf{Titanium Ti-6Al-4V}},
   2.3/{\textbf{Graphite --- CFRP}},
   1.0/{Gold}}
{
  \draw[tick] (\W,\y) -- (\W+0.22,\y);
  \node[mt] at (\W+0.28,\y) {\n};
}

% ---- pairings ----
% (1) titanium in CFRP -- adjacent
\draw[good] (-0.45,3.3) -- (-0.75,3.3) -- (-0.75,2.3) -- (-0.45,2.3);
\node[gl] at (-6.55,2.8) {\textbf{Titanium in CFRP.} Almost no driving voltage --- which is the whole reason Ti-6Al-4V is \emph{the} composite fastener.};
% (2) aluminium in CFRP -- far apart, and the fastener is the small part
\draw[bad] (-1.30,7.4) -- (-1.70,7.4) -- (-1.70,2.3) -- (-1.30,2.3);
\node[bl] at (-7.35,5.1) {\textbf{Aluminium in CFRP.} Far apart \emph{and} the wrong way round: a small aluminium fastener is the \textbf{anode}, the whole carbon panel is the \textbf{cathode}. ``The small anode \dots\ to large cathode \dots\ accelerates the corrosion.''};
% (3) cadmium next to aluminium
\draw[good] (-0.45,8.0) -- (-0.75,8.0) -- (-0.75,7.4) -- (-0.45,7.4);
\node[gl] at (-6.55,8.6) {\textbf{Cadmium-plated steel in aluminium.} Cadmium sits \emph{beside} the aluminium alloys. That is what the plating is for --- not decoration.};

% ---- the two rules ----
\node[font=\footnotesize,align=left,text width=6.6cm,anchor=north west] at (-7.9,0.25)
 {\textbf{Rule 1.} ``The farther apart two materials are in the \dots\ list, the greater the
  galvanic action between them.''\\[3pt]
  \textbf{Rule 2.} ``Because the anode is eroded in a galvanic cell, \emph{it should be the larger
  mass in the cell}.'' A fastener is always the small part --- so the \textbf{fastener must be the
  cathode}.\\[3pt]
  {\scriptsize Both: NASA RP-1228.}};

\end{tikzpicture}
